\section{Detection of primitive error tensors}\label{sec:detection}

The sewing theorem proves scalar identities. We now show that its class of tests detects every coefficient of the Virasoro error, without a restriction on the number of primitive insertions.

\begin{definition}\label{det:setup}
A \emph{detection subspace} for a smooth projective \(n\)-fold \(X\) is a real rational Hodge substructure
\[
 B_0\subset H^*(X;\Q)
\]
such that:
\begin{enumerate}[label=\textup{(\roman*)}]
\item \(B_0\) is spanned by homogeneous coherent admissible singles of type \((j,j)\);
\item the integration pairing is nondegenerate on \(B_0\);
\item \(B_0\) contains all cohomology outside \(H^n(X)\), and contains the nonprimitive middle cohomology for some polarization.
\end{enumerate}
We write \(J=B_0^\perp\), so \(H^*(X;\C)=B_{0,\C}\oplus J_\C\) orthogonally and \(J\subset H^n(X)\) is primitive.
\end{definition}

Only the target on which an error tensor is detected needs a detection subspace. On the shorter targets used in sewing, the same formal combinations of coherent singles need not define projectors.

\begin{lemma}[Available primitive operations]\label{det:operations}
Suppose \(X\) has a detection subspace. Its primitive coevaluation and any function of either Hodge index on a primitive leg are finite linear combinations of the admissible operations of \cref{sec:conventions}. The same holds for the second-index operator on the full cohomology, after decomposition into \(B_0\) and \(J\).
\end{lemma}

\begin{proof}
Choose a homogeneous real rational basis \(b_1,\ldots,b_s\) of \(B_0\) from the coherent singles in \cref{det:setup}, and let \(g^{ab}\) be the inverse of its pairing matrix. Rational linear combinations of those singles remain coherent. Since these classes are even, its coevaluation is
\[
 \Delta_{B_0}=\sum_{a,b}g^{ab}b_a\ot b_b.
\]
Orthogonality and nondegeneracy give
\begin{equation}\label{eq:primitive-coevaluation}
 \Delta_J=\Delta_X-\Delta_{B_0}.
\end{equation}
This is a full coevaluation minus a finite sum of coherent singles. A first-index function on either leg is admissible. On \(J^{p,q}\), one has \(p+q=n\), so a second-index function is the corresponding function of \(n-p\). On each \(b_a\) both indices are known. Decomposing a propagator by \eqref{eq:primitive-coevaluation} therefore realizes the required second-index factors as well.

The same construction gives the projectors onto \(B_0\) and \(J\) by contracting one leg with the integration pairing. All of these are finite diagrams of permitted coevaluations, singles, and cups. On another target, the resulting expressions still make sense as formal test diagrams even if they no longer have the projector interpretation.
\end{proof}

\begin{lemma}[A positive scalar test]\label{det:norm}
Let \(E\in(J_\C^*)^{\ot r}\) be a coefficient tensor obtained by applying one positive Virasoro coefficient rule, with fixed coherent \(B_0\)-inputs in all other slots. Its positive Hermitian squared norm can be written as a finite linear combination of admissible scalar tests of that same coefficient rule.
\end{lemma}

\begin{proof}
The ordinary Gromov--Witten tensors and the positive coefficient operations are even, and every fixed \(B_0\)-input is even. If \(nr\) is odd, the coefficient tensor \(E\) therefore vanishes and the assertion is immediate. In the remaining cases \(E\) is an even tensor.

Let \(Q_J(v,w)=\int_Xv\cup w\) and let \(C_J\) be the Weil operator, acting by \(i^{p-q}\) on \(J^{p,q}\). Hodge--Riemann gives the positive Hermitian form
\begin{equation}\label{eq:primitive-hermitian}
 h_J(v,w)=(-1)^{n(n-1)/2}Q_J(C_Jv,\overline w).
\end{equation}
Our convention is linearity in the first variable. The induced Hermitian form on the dual tensor power is positive definite for every \(r\). In an \(h_J\)-orthonormal basis its squared norm is
\begin{equation}\label{eq:error-norm}
 \|E\|^2=\sum_{a_1,\ldots,a_r}
       \left|E(e_{a_1},\ldots,e_{a_r})\right|^2.
\end{equation}
For \(r=0\) this means the squared absolute value of the scalar error.

More precisely, put \(B_J(v,w)=(-1)^{n(n-1)/2}Q_J(C_Jv,w)\) and let \(K_J=B_J^{-1}\) be its inverse bilinear kernel. Explicitly,
\[
 K_J=(-1)^{n(n-1)/2}(C_J\otimes\id)\Delta_J.
\]
Pair \(E\) with \(E^\#(w_1,\ldots,w_r)=\overline{E(\bar w_1,\ldots,\bar w_r)}\) through one copy of \(K_J\) per corresponding slot. We order the slots as \((1,1'),\ldots,(r,r')\); converting from the block order \((1,\ldots,r,1',\ldots,r')\) contributes the fixed Koszul permutation, which is \((-1)^{r(r-1)/2}\) when \(J\) is odd. Since \(E\) and \(E^\#\) are even, their tensor-product evaluation contributes no further parity sign. Including the interleaving sign in the numerical coefficient makes the contraction equal to \eqref{eq:error-norm}, the ordinary positive tensor norm rather than a supertrace. By \cref{det:operations}, \(K_J\) uses only admissible propagators and first-index functions, since \(q=n-p\) on \(J\). Thus it remains to realize the conjugate coefficient tensor.

Expand \(E\) by \cref{conv:positive-rule}. The ordinary Gromov--Witten correspondence is a real cohomology class: it is obtained from a rational virtual cycle and ordinary evaluation and cotangent operations. The integration pairing and the coherent rational classes are real as well. Complex conjugation consequently replaces each first-index factor by the corresponding second-index factor, conjugates numerical coefficients, and leaves the ordinary Gromov--Witten amplitudes unchanged. A Chern-class multiplication has the fixed bidegree shift \((1,1)\). Use a second ordinary replica, expand the conjugated coefficient rule term by term, and realize its grading factors by \cref{det:operations}. Any classical term remains an allowed cup integral. Extract the specified curve and genus coefficients independently in the two replicas.

There is no assertion that the operator acts on both replicas in the sewing theorem. The first replica carries the positive constraint; every summand of the conjugated expansion on the second replica is part of its fixed test. The latter is a finite diagram at the coefficient in question.

Finally, order all tensor slots once. Every vector in \(J\) has parity \(n\bmod2\). Converting the fixed order in the bilinear diagram to the ordinary positive sum \eqref{eq:error-norm} therefore introduces only the prescribed fixed Koszul signs; these can be included in the numerical coefficients of the tests. In particular, we take the norm in the full underlying tensor power, rather than a supertrace which could cancel positive terms. No bound on \(r\) enters the argument.
\end{proof}

\begin{theorem}[Detection]\label{det:detection}
Assume that \(X\) has a detection subspace, and that all admissible scalar tests of every positive Virasoro coefficient vanish, with degree weights distinguishing the curve classes under consideration. Then \(L_k^XZ_X=0\) for every \(k>0\), for all cohomology insertions and in each such curve class.
\end{theorem}

\begin{proof}
Fix \(k\), genus and curve coefficients, descendant exponents, and a labelled list of slots. Decompose each slot as \(B_{0,\C}\oplus J_\C\). In any summand put specified basis vectors of \(B_0\) into its \(B_0\)-slots. The remaining error is a tensor \(E\) on a power of \(J\). Every scalar test in \cref{det:norm} vanishes by hypothesis; hence \(\|E\|^2=0\). Positivity gives \(E=0\). Since the fixed coefficients and all choices of slots and basis vectors were arbitrary, multilinearity proves the assertion.
\end{proof}

\begin{remark}\label{det:other-targets}
The coefficients in \eqref{eq:primitive-coevaluation} and the functions which describe second Hodge degree are chosen for \(X\). Under a surgery they remain fixed external numerical data in the test. Their failure to describe an orthogonal projection or a positive norm on a shorter target causes no difficulty: the complete Virasoro identity on that target vanishes under every such linear test. Positivity is used only after transport, on \(X\) itself.
\end{remark}
